\documentclass[a4paper,10pt]{article}

\usepackage{marvosym}
\usepackage{courier}
\usepackage{fontspec,url,parskip}
\RequirePackage{color,graphicx}
\usepackage[usenames,dvipsnames]{xcolor}
\usepackage[big]{layaureo}
\usepackage{supertabular}
\usepackage{titlesec}
\usepackage{multirow}
\usepackage{hyperref}
\usepackage{enumitem}
\usepackage{ulem}
\usepackage[absolute]{textpos}
\usepackage{geometry}
\usepackage{tabularx}

\definecolor{linkcolour}{rgb}{0,0.2,0.6}
\hypersetup{
  colorlinks,
  breaklinks,
  urlcolor=linkcolour,
  linkcolor=linkcolour
}

% Load the original icon font by filename, without requiring a system font install.
\newfontfamily{\ResumeLinkFont}{FontAwesome.otf}
\newcommand{\faPaperclip}{{\ResumeLinkFont\symbol{"F0C6}}}

\defaultfontfeatures{Mapping=tex-text}
\titleformat{\section}{\Large\scshape\raggedright}{}{0em}{}[\titlerule]
\titlespacing{\section}{0pt}{3pt}{3pt}
\hyphenation{im-pre-se}

\setlength{\TPHorizModule}{30mm}
\setlength{\TPVertModule}{\TPHorizModule}
\textblockorigin{2mm}{0.65\paperheight}
\setlength{\parindent}{0pt}

\geometry{
  a4paper,
  total={200mm,257mm},
  left=0.8in,
  right=0.8in,
  top=0.5in,
  bottom=0.7in
}

\begin{document}

% GENERATED:ACADEMIC_PROFILE:START
\par{\centerline{
  {\Huge George Stepaniants}}\bigskip
  \centerline{\href{\detokenize{mailto:gstepan@caltech.edu}}{gstepan@caltech.edu} \,|\, \href{\detokenize{mailto:georgestepaniants@gmail.com}}{georgestepaniants@gmail.com} \,|\, \href{\detokenize{https://georgestepaniants.com}}{https://georgestepaniants.com}}
}\bigskip

I develop data-driven methods to learn laws and relationships in scientific data, using foundations of machine learning, statistical theory, numerical analysis, and mathematical modeling to design approaches that succeed in data-limited regimes and embody domain-specific inductive biases. My teaching philosophy is inspired by my research, showing students how to discover mathematical ideas in field-specific literature, translate them into well-posed theories, and bring these theories to life as numerical algorithms and reproducible code.

\underline{\smash{Research Areas:}} Applied and Computational Mathematics, Statistical Theory and Methods, High-Dimensional Statistics, Machine Learning, Scientific Computing, Numerical Analysis, Scientific ML, Mechanics of Materials, Fluid Dynamics, Biophysics

\section{Academic Employment}
\begin{tabular}{lp{13cm}}
(To Begin) \textsc{Oct} 2026 & \textbf{University of Cambridge (Assistant Professor)}\\
& \textit{Assistant Professor in Mathematics of Information}\\
& Department of Applied Mathematics and Theoretical Physics (DAMTP)\\
\end{tabular}\\

\begin{tabular}{lp{13cm}}
(To Begin) \textsc{Oct} 2026 & \textbf{New York University (Visiting Faculty)}\\
& Courant Institute School of Mathematics, Computing, and Data Science\\
\end{tabular}\\

\begin{tabular}{lp{13cm}}
\textsc{Sep} 2024 -- \textsc{Present} & \textbf{California Institute of Technology (Postdoctoral Scholar)}\\
& \textit{NSF Mathematical Sciences Postdoctoral Research Fellow (MSPRF)}\\
& Department of Computing and Mathematical Sciences (CMS)\\
& Postdoctoral Advisor: \href{\detokenize{http://stuart.caltech.edu/}}{Andrew M. Stuart}\\
\end{tabular}\\

\section{Education}
\begin{tabular}{lp{13cm}}
\textsc{Sep} 2019 -- \textsc{Jun} 2024 & \textbf{Massachusetts Institute of Technology (PhD)}\\
\textsc{\bf GPA}: 4.9/5.0 & Department of Mathematics and Institute for Data, Systems, and Society (IDSS)\\
& PhD Advisors: \href{\detokenize{https://math.mit.edu/~rigollet/}}{Philippe Rigollet} and \href{\detokenize{https://math.mit.edu/~dunkel/}}{Jörn Dunkel}\\
& Graduate Thesis: \textit{\href{\detokenize{https://dspace.mit.edu/handle/1721.1/155887}}{Inference from limited observations in statistical, dynamical, and functional problems}}\\
\\
\textsc{Sep} 2015 -- \textsc{Jun} 2019 & \textbf{University of Washington (BSc)}\\
\textsc{\bf GPA}: 3.87/4.00 & Department of Mathematics and Department of Computer Science (double major)\\
& Undergraduate Research Advisors: \href{\detokenize{http://faculty.washington.edu/kutz/}}{Nathan Kutz} and \href{\detokenize{https://www.bingbrunton.com/}}{Bing Brunton}\\
& Research Topic: \textit{\href{\detokenize{https://journals.aps.org/pre/abstract/10.1103/PhysRevE.102.042309}}{Inferring causal networks of dynamical systems through transient dynamics and perturbation}}\\
\\
\end{tabular}\\

\section{Academic Awards}
\renewcommand{\arraystretch}{1.3}
\begin{tabularx}{\textwidth}{X|r}
Isaac Newton Institute 2026 Simons Fellow & \textsc{Aug} 2026\\
ILAS 2026 Early Career NSF Conference Grant Award Recipient & \textsc{Apr} 2026\\
NSF Mathematical Sciences Postdoctoral Research Fellowship (MSPRF) \href{\detokenize{https://www.nsf.gov/awardsearch/showAward?AWD_ID=2402074&HistoricalAwards=false}}{2402074} & \textsc{Sep} 2024 -- \textsc{Current}\\
IMS Lawrence D. Brown Ph.D. Student Award Recipient \href{\detokenize{https://imstat.org/wp-content/uploads/2025/02/Bulletin54_3-embargo.pdf}}{\faPaperclip} (3 recipients nationwide) & \textsc{Sep} 2024\\
NSF Graduate Research Fellowship (GRFP) \href{\detokenize{https://www.nsf.gov/awardsearch/showAward?AWD_ID=1745302}}{1745302} & \textsc{Jun} 2019 -- \textsc{Jun} 2024\\
SIAM Student Travel Award & \textsc{Dec} 2023\\
Calouste Gulbenkian Foundation Short Term Conference and Travel Grant & \textsc{Jun} 2023\\
MIT Presidential Fellowship & \textsc{Sep} 2019 -- \textsc{Jun} 2020\\
Phi Beta Kappa Honors Society Member & \textsc{Jun} 2019\\
Mary Gates Research Scholarship (merit-based) & \textsc{Jun} 2019\\
University of Washington Dean's List & \textsc{Sep} 2015 -- \textsc{Jun} 2019\\
Early acceptance to University of Washington (UW Academy) & \textsc{Sep} 2015\\
\end{tabularx}
\vspace{10pt}

\section{Publications}
Please see also my \href{\detokenize{https://scholar.google.com/citations?user=CKYZLxYAAAAJ&hl=en}}{Google Scholar} and \href{\detokenize{https://orcid.org/0000-0002-7834-7536}}{ORCID ID}. ($\dagger$ alphabetical order, $^*$ joint first author)
\vspace{-7pt}
\subsubsection*{Thesis}
\begin{enumerate}[label={[T\arabic*]}, leftmargin=*]
\item \href{\detokenize{https://dspace.mit.edu/handle/1721.1/155887?show=full}}{\faPaperclip}~ George Stepaniants. ``Inference from Limited Observations in Statistical, Dynamical, and Functional Problems.'' \textit{Massachusetts Institute of Technology} (2024).
\end{enumerate}

\subsubsection*{Preprints}
\begin{enumerate}[label={[PP\arabic*]}, leftmargin=*]
\item \href{\detokenize{https://doi.org/10.5281/zenodo.22289731}}{\faPaperclip}~ Matthew J. Colbrook, George Stepaniants, and Alex Townsend. ``A Complete Resolution of Forsythe's Conjecture for Restarted Conjugate Gradients.'' \textit{Zenodo preprint} (2026). DOI: 10.5281/zenodo.22289731.

\item \href{\detokenize{https://arxiv.org/abs/2608.08953}}{\faPaperclip}~ Matthew J. Colbrook, Siavash Sadeghi, and George Stepaniants. ``A computer-assisted counterexample to the planar Berenstein conjecture.'' \textit{arXiv preprint arXiv:2608.08953} (2026).

\item \href{\detokenize{https://arxiv.org/abs/2608.02852}}{\faPaperclip}~ Matthew J. Colbrook, George Stepaniants, and Alex Townsend. ``A Proof of the Forsythe Conjecture for the Two-Step Restarted Conjugate Gradient Method.'' \textit{arXiv preprint arXiv:2608.02852} (2026).

\item \href{\detokenize{https://arxiv.org/abs/2608.01579}}{\faPaperclip}~ Matthew J. Colbrook and George Stepaniants. ``A computer-assisted counterexample to the planar Pompeiu and Schiffer conjectures.'' \textit{arXiv preprint arXiv:2608.01579} (2026).
\end{enumerate}

\subsubsection*{Journal Articles}
\begin{enumerate}[label={[J\arabic*]}, leftmargin=*]
\item \href{\detokenize{https://smai-jcm.centre-mersenne.org/articles/10.5802/smai-jcm.148/}}{\faPaperclip}~ Kaushik Bhattacharya, Lianghao Cao, $\dagger$ George Stepaniants, Andrew M. Stuart, and Margaret Trautner. ``Learning Memory and Material Dependent Constitutive Laws.'' \textit{The SMAI Journal of Computational Mathematics} 12 (2026): 219-267. $\dagger$ Authors listed alphabetically.

\item \href{\detokenize{https://doi.org/10.1090/cams/64}}{\faPaperclip}~ David Darrow and $\dagger$ George Stepaniants. ``A Spectral Theory of Scalar Volterra Equations.'' \textit{Communications of the American Mathematical Society} 6 (2026): 634-725. $\dagger$ Authors listed alphabetically.

\item \href{\detokenize{https://epubs.siam.org/doi/full/10.1137/24M1682841}}{\faPaperclip}~ Yanjun Han, Philippe Rigollet, and $\dagger$ George Stepaniants. ``Covariance alignment: from maximum likelihood estimation to Gromov-Wasserstein.'' \textit{SIAM Journal on Mathematics of Data Science} 7.3 (2025): 1491-1513. $\dagger$ Authors listed alphabetically.

\item \href{\detokenize{https://journals.aps.org/prresearch/abstract/10.1103/PhysRevResearch.6.043062}}{\faPaperclip}~ George Stepaniants$^*$, Alasdair D. Hastewell$^*$, Dominic J. Skinner, Jan F. Totz, and Jörn Dunkel. ``Discovering dynamics and parameters of nonlinear oscillatory and chaotic systems from partial observations.'' \textit{Physical Review Research} 6, 043062 (2024). $^*$ Equal contribution.

\item \href{\detokenize{https://elifesciences.org/articles/91597}}{\faPaperclip}~ Marie Breeur$^*$, George Stepaniants$^*$, Pekka Keski-Rahkonen, Philippe Rigollet, and Vivian Viallon. ``Optimal transport for automatic alignment of untargeted metabolomic data.'' \textit{eLife} 12:RP91597 (2024). $^*$ Equal contribution.

\item \href{\detokenize{https://jmlr.org/papers/v24/21-1363.html}}{\faPaperclip}~ George Stepaniants. ``Learning partial differential equations in reproducing kernel Hilbert spaces.'' \textit{Journal of Machine Learning Research} 24.86 (2023): 1-72.

\item \href{\detokenize{https://journals.aps.org/pre/abstract/10.1103/PhysRevE.102.042309}}{\faPaperclip}~ George Stepaniants, Bingni W. Brunton, and J. Nathan Kutz. ``Inferring causal networks of dynamical systems through transient dynamics and perturbation.'' \textit{Physical Review E} 102.4 (2020): 042309.
\end{enumerate}

\subsubsection*{Conference Proceedings}
\begin{enumerate}[label={[C\arabic*]}, leftmargin=*]
\item \href{\detokenize{https://arxiv.org/abs/2210.06545}}{\faPaperclip}~ Enric Boix-Adserà, Hannah Lawrence, George Stepaniants, and Philippe Rigollet. ``GULP: a prediction-based metric between representations.'' \textit{Advances in Neural Information Processing Systems} (2022).

\item \href{\detokenize{https://proceedings.mlr.press/v130/chewi21a}}{\faPaperclip}~ Sinho Chewi, Julien Clancy, Thibaut Le Gouic, Philippe Rigollet, $\dagger$ George Stepaniants, and Austin Stromme. ``Fast and smooth interpolation on Wasserstein space.'' \textit{International Conference on Artificial Intelligence and Statistics} PMLR (2021). $\dagger$ Authors listed alphabetically.
\end{enumerate}

\subsubsection*{In Preparation}
\begin{enumerate}[label={[P\arabic*]}, leftmargin=*]
\item David Darrow$^*$, George Stepaniants$^*$, and Chris Camaño. ``SIEVE: Spectral integral transforms, poly-exponential approximants, and Volterra equations.'' $^*$ Equal contribution.
\end{enumerate}
\vspace{10pt}

\section{Talks and Presentations}
\subsubsection*{Organized Symposia}
\begin{itemize}
\item ``Minisymposium on Data-Driven Methods for Multiscale Modeling and Homogenization'', SIAM Conference on Computational Science and Engineering, Fort Worth, \textsc{March} 2025

\item ``Minisymposium on Data-Driven Learning of Dynamical Systems from Partial Observations'', SIAM Conference on Mathematics of Data Science, Atlanta, \textsc{October} 2024
\end{itemize}

\subsubsection*{Invited Talks}
\begin{itemize}
\item ``Volterra Integral Equations and Memory Dependent Constitutive Laws'', Isaac Newton Institute, Cambridge, \textsc{August} 2026

\item ``Learning Material Constitutive Laws with Neural Operators'', ILAS, Virginia Tech, \textsc{May} 2026

\item ``Volterra Integral Equations and Memory Dependent Constitutive Laws'', UC Irvine Applied \& Computational Math Seminar, Irvine, \textsc{October} 2025

\item ``Learning Memory and Material Dependent Constitutive Laws'', Surrogates and Dimension Reduction in Scientific Machine Learning, Manchester University, \textsc{September} 2025

\item ``Alignment of Untargeted Data through their Covariances: A Novel Perspective on a Classical Tool in Optimal Transport'', Joint Statistics Meeting, Nashville, \textsc{August} 2025 \textbf{(One of 3 PhD students selected for the prestigious IMS Lawrence D. Brown PhD Student Award)} \href{\detokenize{https://imstat.org/2025/04/02/preview-george-stepaniants-ims-lawrence-d-brown-phd-student-lecture/}}{\faPaperclip}

\item ``A Spectral Theory of Volterra Equations: Applications to Learning of Material Laws'', Efficient and Reliable Deep Learning Methods and their Scientific Applications, \textbf{Banff BIRS Centre}, \textsc{June} 2025

\item ``Learning Dynamics of Hidden Variables in Multiscale Viscoelastic Materials'', SIAM Conference on Applications of Dynamical Systems, Denver, \textsc{May} 2025

\item ``A Spectral Theory of Scalar Volterra Equations'', Dartmouth Applied \& Computational Math Seminar, Dartmouth, \textsc{March} 2025

\item ``A Spectral Theory of Scalar Volterra Equations'', MIT Applied Math Physical Mathematics Seminar, Cambridge, \textsc{March} 2025

\item ``Learning Memory and Material Dependent Constitutive Laws'', Differential Equations for Data Science, Kyoto University, \textsc{February} 2025

\item ``Discovering dynamics and parameters of nonlinear oscillatory and chaotic systems from partial observations'', Fourth Symposium on Machine Learning and Dynamical Systems, \textbf{Fields Institute}, \textsc{July} 2024

\item ``Covariance Alignment with Optimal Transport'', Yale Applied Mathematics Seminar, New Haven, \textsc{April} 2024

\item ``Gromov-Wasserstein Theory and Application to Metabolomics'', SIAM Conference on Uncertainty Quantification, Trieste, Italy, \textsc{February} 2024

\item ``Gromov-Wasserstein Theory and Application to Metabolomics'', Statistics and Learning Theory Summer School, Tsaghkadzor, Armenia, \textsc{July} 2023 \textbf{(One of 7 invited speakers)}

\item ``Optimal transport for automatic alignment of untargeted metabolomic data'', Harvard Applied Math Graduate Student Seminar, Cambridge, \textsc{March} 2023

\item ``Learning PDEs in a Reproducing Kernel Hilbert Space'', SIAM Conference on Mathematics of Data Science, San Diego, \textsc{September} 2022

\item ``Learning PDEs in a Reproducing Kernel Hilbert Space'', Meeting on Mathematical Statistics, \textbf{CIRM, Marseille, France}, \textsc{December} 2021 \textbf{(Only 3 graduate student speakers invited)}
\end{itemize}

\subsubsection*{Contributed Talks}
\begin{itemize}
\item ``Discovering dynamics and parameters of nonlinear oscillatory and chaotic systems from partial observations'', Dynamics Days, UC Davis, \textsc{January} 2024

\item ``Learning and predicting complex systems dynamics from single-variable observations'', APS March Meeting, Chicago, \textsc{March} 2022

\item ``Learning PDEs in a Reproducing Kernel Hilbert Space'', LIDS Stats \& Tea, MIT, \textsc{December} 2021

\item ``Inferring causal networks of dynamical systems through transient dynamics and perturbation'', Econometrics Lunch, MIT, \textsc{December} 2021

\item ``Fusion of Genetically Incompatible Fungal Cells'', UCLA Computational and Applied Math REU Presentation, IPAM, \textsc{August} 2018

\item ``Quantifying Rupture Risk of Brain Aneurysms'', MATDAT18: NSF Materials and Data Science Hackathon, Alexandria, \textsc{June} 2018 \href{\detokenize{https://matdat18.wordpress.ncsu.edu/files/2018/06/Team12.pdf}}{\faPaperclip}

\item ``Hyperparameter Selection'', AI2 Research Internship Final Presentation, Seattle, \textsc{August} 2017

\item ``Beaker Experimentation Platform'', AI2 Research Internship Midterm Presentation, Seattle, \textsc{August} 2017

\item ``Image Analysis in Parkinson's Research'', Pfizer Research Internship Final Presentation, Cambridge, \textsc{August} 2016
\end{itemize}

\subsubsection*{Poster Presentations}
\begin{itemize}
\item ``Covariance alignment: from maximum-likelihood estimation to Gromov-Wasserstein'', Cornell ORIE Young Researchers Workshop, Cornell, \textsc{October} 2023

\item ``Inferring causal networks of dynamical systems through transient dynamics and perturbation'', Undergraduate Research Symposium, UW, \textsc{June} 2019
\end{itemize}

\section*{Teaching Experience}
\begin{tabularx}{\textwidth}{@{}X|>{\raggedleft\arraybackslash}p{2.2cm}@{}}
\textbf{(ACM 270) California Institute of Technology -- Instructor of Record} \href{\detokenize{https://schedules.caltech.edu/syllabi/ACM270_SP2425.pdf}}{\faPaperclip} & \textsc{Spring} 2025\\[2pt]
Designer and Instructor of New Course for Data-Driven Modeling of Dynamical Systems\\
\textit{Two weekly 1.5 hour classes with 30+ students, with final presentations and projects, unifying topics in system identification, sparse model inference, Koopman theory, adjoint methods, data assimilation, and scientific ML. Turning class content and lectures into a class text.}
\end{tabularx}
\vspace{10pt}

\begin{tabularx}{\textwidth}{@{}X|>{\raggedleft\arraybackslash}p{2.2cm}@{}}
\textbf{(MIT 18.032) Massachusetts Institute of Technology -- Teaching Assistant} \href{\detokenize{http://student.mit.edu/catalog/search.cgi?search=18.032&style=verbatim}}{\faPaperclip} & \textsc{Spring} 2022\\[2pt]
Teaching Assistant for Differential Equations (theory focused)\\
\textit{Two weekly 1 hour recitations with 16 students, including grading of homeworks/finals and proctoring final exams. Covering existence and uniqueness theory (Cauchy-Lipschitz Theorem, Picard iterations), control and comparison of solutions, maximal solutions, Grönwall's Lemma, multivariable ODEs, and fixed point stability.}
\end{tabularx}
\vspace{10pt}

\begin{tabularx}{\textwidth}{@{}X|>{\raggedleft\arraybackslash}p{2.2cm}@{}}
\textbf{(MIT 18.600) Massachusetts Institute of Technology -- Teaching Assistant} \href{\detokenize{http://student.mit.edu/catalog/search.cgi?search=18.600&style=verbatim}}{\faPaperclip} & \textsc{Fall} 2021\\[2pt]
Teaching Assistant for Introduction to Probability\\
\textit{Two weekly 1 hour recitations with 40 students, including grading of homeworks and finals. Covering probability spaces, random variables, classes of distributions, Bayes theorem, Chebyshev inequality, law of large numbers, and central limit theorems.}
\end{tabularx}
\vspace{10pt}

\section*{Mentoring Experience}
\begin{tabularx}{\textwidth}{@{}X|>{\raggedleft\arraybackslash}p{2.2cm}@{}}
\textbf{(SURF) California Institute of Technology} & \textsc{Summer} 2025\\[2pt]
Undergraduate Research Mentor for Vladislav Syntko (Maastricht University, Netherlands)\\
A Signature-Based Approach for System Identification and Control: \textit{Applications and theory for signature transform methods in open-loop control of dynamical systems.}
\end{tabularx}
\vspace{10pt}

\begin{tabularx}{\textwidth}{@{}X|>{\raggedleft\arraybackslash}p{2.2cm}@{}}
\textbf{(WAVE) California Institute of Technology} & \textsc{Summer} 2025\\[2pt]
Undergraduate Research Mentor for Owen Tolbert (University of Maryland Baltimore County)\\
A Study of Network Inference Methods: \textit{Information-theoretic and deep learning methods for inference of networked dynamical systems.}
\end{tabularx}
\vspace{10pt}

\begin{tabularx}{\textwidth}{@{}X|>{\raggedleft\arraybackslash}p{2.2cm}@{}}
\textbf{(MCM) California Institute of Technology} & \textsc{Fall} 2024\\[2pt]
Trained three undergraduate Caltech teams for the Mathematical Contest in Modeling (MCM), coaching and solving practice problems over the course of several months
\begin{itemize}
  \item (Honorable Mention) Gautham Kappaganthula, Constantin Cedillo-Vayson de Pradenne, Colin La
  \item (Successful Participant) Joseph Pieper, Sujay Champati, Dhruv Verma
  \item (Successful Participant) James Hou, Aman Burman, Abhiram Cherukupalli
\end{itemize}
\end{tabularx}
\vspace{10pt}

\begin{tabularx}{\textwidth}{@{}X|>{\raggedleft\arraybackslash}p{2.2cm}@{}}
\textbf{(Mentor) California Institute of Technology} & \textsc{Fall} 2024\\[2pt]
Undergraduate Research Mentor for Zixiang Zhou (University of Southern California)\\
\textit{Mentored research reading and project in data-driven dynamical systems inference algorithms based on the method of characteristics.}
\end{tabularx}
\vspace{10pt}

\begin{tabularx}{\textwidth}{@{}X|>{\raggedleft\arraybackslash}p{2.2cm}@{}}
\textbf{(SPUR+) Massachusetts Institute of Technology} & \textsc{Summer} 2023 -- \textsc{Fall} 2023\\[2pt]
Undergraduate Research Mentor for Elaine Liu (Massachusetts Institute of Technology)\\
\href{\detokenize{https://math.mit.edu/documents/spur/2023/Liu.pdf}}{Modeling International Trade and Tariffs}: \textit{Study of large trade and tariffs dataset across 200 world countries, investigating the use of spectral and graph wavelet decompositions for analysis of temporal trade network data.}
\end{tabularx}
\vspace{10pt}

\begin{tabularx}{\textwidth}{@{}X|>{\raggedleft\arraybackslash}p{2.2cm}@{}}
\textbf{(DRP) Massachusetts Institute of Technology} & \textsc{Summer} 2023\\[2pt]
Directed Reading Program with Loreta Arzumanyan and Joshua Curtis Kuffour (Massachusetts Institute of Technology)\\
\textit{Guided reading of two undergraduate students in the graduate dynamical systems text "Stability, Instability and Chaos" by Paul Glendinning over the course of the summer. Prepared students to present their knowledge of the text in a final presentation at the end of summer.}
\end{tabularx}
\vspace{10pt}

\begin{tabularx}{\textwidth}{@{}X|>{\raggedleft\arraybackslash}p{2.2cm}@{}}
\textbf{(Mentor) Massachusetts Institute of Technology} & \textsc{Summer} 2023\\[2pt]
Undergraduate Research Mentor for Donald J. Liveoak and Hanna Chen (Massachusetts Institute of Technology)\\
\textit{Guided research readings with two MIT undergraduates on optimal transport and adjoint methods for inference of stochastic dynamical systems and networked dynamical systems.}
\end{tabularx}
\vspace{10pt}

\begin{tabularx}{\textwidth}{@{}X|>{\raggedleft\arraybackslash}p{2.2cm}@{}}
\textbf{(UROP) Massachusetts Institute of Technology} & \textsc{Fall} 2021 -- \textsc{Spring} 2022\\[2pt]
Undergraduate Research Mentor for David Darrow (Massachusetts Institute of Technology)\\
Optimal Transport for Protein Folding: \textit{Studying how optimal transport and Gromov-Wasserstein methods can be used to predict the three-dimensional structure of proteins.}
\\
\underline{\smash{David Darrow awarded 2022 Churchill Scholarship}} \href{\detokenize{https://news.mit.edu/2022/three-with-mit-ties-win-churchill-scholarships-0112}}{\faPaperclip}
\end{tabularx}
\vspace{10pt}

\begin{tabularx}{\textwidth}{@{}X|>{\raggedleft\arraybackslash}p{2.2cm}@{}}
\textbf{(Community Service) University of Washington}\\
Math tutoring from K12 to college-level subjects & 2015 -- 2019\\
Teaching assistant at University of Washington Math Circle & 2015 -- 2016\\
\end{tabularx}
\vspace{10pt}

\section{Service and Leadership}
\begin{enumerate}
  \item Organizer of Caltech CMS Departmental CMX Seminar 2025 -- 2026
  \item Board Member of One World Seminar on Mathematics of Machine Learning 2025 \href{\detokenize{https://sites.google.com/view/oneworldml}}{\faPaperclip}
  \item Organizer of SURF/WAVE undergraduate research across three faculty at Caltech CMS 2025
  \item Presenter on AI literacy and societal impact in LA and NY Armenian communities 2025
  \item Organized panel on graduate school experience and research in Redmond Armenian community 2020
  \item Founded and led Armenian Student Association at the University of Washington (ASAUW) 2015 -- 2019
  \item Competition judge at the University of Washington Math Olympiad 2015 -- 2019
\end{enumerate}
\vspace{10pt}

\section{Professional Membership}
\begin{itemize}
  \item Institute of Mathematical Statistics (IMS) 2024 -- \textsc{Present}
  \item Society for Industrial and Applied Mathematics (SIAM) 2020 -- \textsc{Present}
  \item American Physical Society (APS) 2019 -- \textsc{Present}
\end{itemize}
\vspace{10pt}

\section{Internship and Research Experience}
\begin{tabular}{r|p{13cm}}
\textsc{Jul} 2018 -- \textsc{Aug} 2018 & Computational and Applied Math Research Experience for Undergrads (REU) at \textbf{UCLA}\\
& \emph{Undergraduate Researcher in Mycofluidics Lab}\\
& \footnotesize{Worked in Professor Marcus Roper's lab on imaging of fungal cells in Neurospora and Ashbya fungal species. Collected data on nuclear division of these multinucleated cells using 3D imaging algorithms and performed data analysis on nuclear spacing and mixing within the cell. Discovered that genetically incompatible fungal strains fuse together when exposed to environmental stress. We are working on a paper that relies on the research findings and imaging algorithms I developed at UCLA.}\\
\multicolumn{2}{c}{}\\
\textsc{Jun} 2017 -- \textsc{Sep} 2017 & Engineering Intern at \textbf{Allen Institute for Artificial Intelligence (AI2)}\\
& \emph{Full Stack Development and Data Analysis/Visualization}\\
& \footnotesize{Collaborated with researchers and built a system which optimized hyperparameter selection in various neural network experiments run in the company. I was responsible for experiment design decisions and execution of these experiments on a Google cluster. My platform significantly simplified the experimentation process and reduced runtime by a factor of two.}\\
\multicolumn{2}{c}{}\\
\textsc{Aug} 2016 -- \textsc{Dec} 2016 & Natural Language Processing Intern at \textbf{ABBYY}\\
& \emph{Parser Accuracy Scoring}\\
& \footnotesize{Compared ABBYY's parser efficiency and output to that of MaltParser. Created scripts in Java to read parser output from CoNLL-X data files and scored them using unlabeled and labeled attachment scores (UAS and LAS). Used Excel for visualization and graphing.}\\
\multicolumn{2}{c}{}\\
\textsc{Jun} 2016 -- \textsc{Aug} 2016 & Image Analysis Intern at \textbf{Pfizer}\\
& \emph{Imaging Algorithms for Automated Brain Slice Imaging}\\
& \footnotesize{Worked in the Neuroscience Pain Research Unit at Pfizer. Studied how various drugs help regenerate healthy cells damaged by neurodegenerative diseases, especially Parkinson's. Used image analysis algorithms, particularly 3D Watershed Segmentation, to quantify the percentage of regenerated healthy cells after treatment. My automated imaging pipeline was used by researchers to quantify hundreds of drug profiles and reduced the runtime of their image analysis code by 10 times.}\\
\end{tabular}
% GENERATED:ACADEMIC_PROFILE:END

\end{document}
